\section{Minimizing velocities and the branch Hessian}
\label{sec:preliminaries}

Throughout the remaining sections we assume the weak MTW condition
\eqref{eq:weak-mtw}. A velocity is called \emph{nonconjugate} if the
vertical differential $d_w\exp_x:T_xM\to T_{\exp_x w}M$ is nonsingular.
We write
\begin{align*}
 L_x(q)&=\{w\in T_xM:\exp_x w=q,\ |w|_x=d(x,q)\},\\
 G_x&=\{w\in T_xM:w\in L_x(\exp_x w)\},\\
 \mathcal C&=\{(x,w)\in TM:w\in G_x\},
 \qquad
 \mathcal I=\{(x,w)\in TM:w\in I(x)\},
\end{align*}
and set $D=\diam(M)$.  Existence of minimizing geodesics shows that
$L_x(q)$ is nonempty.  Every velocity in $\mathcal C$ has norm at most
$D$.  Continuity of the distance and the exponential map shows that
$\mathcal C$ is closed; as the base manifold is compact, $\mathcal C$ is
therefore compact.

We first record the index-form facts used below. For geometric background,
see \cite[Sections~5.7--5.8]{Calegari2013}; for the unrestricted energy
second-variation formula, see \cite[Proposition~5.3.8]{Gorodski2016}.
The latter is the revised course chapter, whose formula allows tangential
as well as normal variation fields. Along a geodesic
$\gamma:[0,T]\to M$, let
\begin{equation}
 Q_\gamma(X,Y)=\int_0^T
 \bigl(\langle D_\tau X,D_\tau Y\rangle
       -\langle R(X,\dot\gamma)\dot\gamma,Y\rangle\bigr)\,d\tau .
 \label{eq:index-form}
\end{equation}
Our curvature convention gives the Jacobi equation
$D_\tau^2X+R(X,\dot\gamma)\dot\gamma=0$.
Fields in \eqref{eq:index-form} may be continuous and piecewise smooth.
A constant-speed minimizing geodesic on $[0,T]$, with $T>0$, also
minimizes the fixed-endpoint energy
$E(\alpha)=\tfrac12\int_0^T|\dot\alpha|^2\,d\tau$.  Indeed,
Cauchy--Schwarz gives
\[
 E(\alpha)\geq\frac{\operatorname{length}(\alpha)^2}{2T}
 \geq\frac{d(\gamma(0),\gamma(T))^2}{2T}=E(\gamma).
\]
Energy second variation therefore gives $Q_\gamma(Z,Z)\geq0$ for every
field with $Z(0)=Z(T)=0$, without a normality restriction.
Approximation gives the same assertion for piecewise smooth fields.  Moreover, equality forces
$Z$ to be a Jacobi field.  Indeed, nonnegativity applied to $Z+tW$
implies $Q_\gamma(Z,W)=0$ for every endpoint-zero test field $W$.
Integration by parts first gives the Jacobi equation on each smooth
piece and then continuity of $D_\tau Z$ at the break points.  Thus $Z$
is a Jacobi field on the entire interval.

\begin{lemma}[Interior characterization]
\label{lem:interior-characterization}
For every $x\in M$,
\begin{equation}
 I(x)\subset G_x,\qquad rG_x\subset I(x)\quad(0\leq r<1).
 \label{eq:radial-interior}
\end{equation}
A velocity $w\in T_xM$ belongs to $I(x)$ if and only if it is the
unique minimizing velocity from $x$ to $\exp_x w$ and is nonconjugate.
The set $\mathcal I$ is open, and $c$ is smooth in a neighborhood of
each pair $(x,\exp_x w)$ with $w\in I(x)$.  There is a number
$\rho>0$ such that $|w|_x<\rho$ implies $w\in I(x)$ for every $x$.
\end{lemma}

\begin{proof}
The first inclusion in \eqref{eq:radial-interior} follows from the
definition.  For $0<r<1$, a proper initial portion of the minimizing
geodesic with velocity $w\in G_x$ remains minimizing up to time
$1/r>1$ when parametrized with initial velocity $rw$.  Hence
$rw\in I(x)$.  The case $r=0$ follows from the constant geodesic.

Suppose $w\in I(x)$, and put $\gamma(\tau)=\exp_x(\tau w)$.
There is a $T>1$ for which $\gamma|_{[0,T]}$ is minimizing.
Any other minimizing geodesic from $x$ to $\gamma(1)$, followed by
$\gamma|_{[1,T]}$, would be a minimizing curve of the same length
from $x$ to $\gamma(T)$.  A minimizing curve has no corner, so its
velocities agree at the joining point.  Uniqueness for the geodesic
equation then identifies the two initial geodesics.
If $d_w\exp_x$ were singular, a nonzero kernel vector would give a
nonzero Jacobi field on $[0,1]$ vanishing at both endpoints.
Extend this field by zero to $[0,T]$.  Its index form is zero, so the
preceding observation makes the extension a Jacobi field on $[0,T]$.
Its vanishing on $(1,T)$ contradicts uniqueness for the Jacobi equation.
Thus $w$ is nonconjugate.

Conversely, suppose that $w$ is a unique minimizing velocity and is
nonconjugate.  The map
\begin{equation}
 \Psi:TM\longrightarrow M\times M,
 \qquad \Psi(z,u)=(z,\exp_z u)
 \label{eq:exponential-pair-map}
\end{equation}
has an invertible differential at $(x,w)$.  Choose neighborhoods on
which it is a diffeomorphism.  After shrinking its image neighborhood,
every minimizing velocity for every pair in that neighborhood belongs
to this inverse branch.  To see this, a sequence of contrary minimizing
velocities would have, by compactness of $\mathcal C$, a subsequential
limit minimizing from $x$ to $\exp_x w$.  Uniqueness makes this limit
$w$, contradicting its exclusion from the open branch neighborhood.
Every nearby pair has a minimizer, so its inverse-branch velocity is
minimizing and is the unique minimizer.  Thus an open neighborhood of
$(x,w)$ in $TM$ consists of minimizing velocities.  A small forward
radial change of $w$ within this neighborhood proves $w\in I(x)$.
For $w=0$, the same conclusion follows directly from the definition;
also $L_x(x)=\{0\}$ and $d_0\exp_x=\Id$.

The inverse neighborhoods just constructed show that $\mathcal I$ is
open.  On their image neighborhoods the true cost is half the squared
norm of the smooth inverse-branch velocity, so it is smooth.
Finally, $\mathcal I$ contains the zero section.  Its openness and
compactness of the zero section give a uniform $\rho>0$ as asserted.
\end{proof}

\subsection{Smooth geodesic branches}

At any nonconjugate $(x,w)$, the local inverse of
\eqref{eq:exponential-pair-map} defines a smooth branch cost
$\hat c(z,q)$: it is half the squared norm of the inverse-branch
velocity.  Define the symmetric bilinear form
\begin{equation}
 A_x(w)=\Hess_{xx}\hat c(x,\exp_x w),
 \label{eq:branch-hessian}
\end{equation}
where the second endpoint is held fixed in taking the covariant
Hessian.  Local inverse branches agree near a prescribed velocity,
so $A$ depends smoothly on $(x,w)$ throughout the nonconjugate set.
When $w\in I(x)$, the branch cost agrees locally with $c$ by
\cref{lem:interior-characterization}.  At a nonconjugate minimizing
cut velocity its central value is still $c(x,\exp_x w)$, and
\eqref{eq:branch-hessian} defines its smooth geodesic-branch Hessian.
In particular, as velocities in $\mathcal I$ approach a nonconjugate
minimizing velocity, their true-cost Hessians converge to this branch
Hessian.  This assertion also holds with the base point varying, using
local tangent-bundle coordinates.

First variation gives $\nabla_x\hat c(x,\exp_x w)=-w$.
For small $\tau$, with $\exp_x w$ fixed, the branch velocity at
$\gamma_{x,w}(\tau)=\exp_x(\tau w)$ is
$(1-\tau)\dot\gamma_{x,w}(\tau)$.  Its negative has covariant
derivative $w$ at $\tau=0$.  Consequently,
\begin{equation}
 A_x(w)(w,\cdot)=\langle w,\cdot\rangle_x.
 \label{eq:radial-hessian}
\end{equation}
This identity is valid at every nonconjugate velocity.

\begin{lemma}[Index-form representation]
\label{lem:hessian-index}
If $w\in G_x$ is nonconjugate, then, along
$\gamma_{x,w}(\tau)=\exp_x(\tau w)$ on $[0,1]$,
\begin{equation}
 A_x(w)(\xi,\xi)
 =\min_{\substack{X(0)=\xi\\X(1)=0}} Q_{\gamma_{x,w}}(X,X).
 \label{eq:hessian-index-minimum}
\end{equation}
The minimum is attained uniquely by the Jacobi field with these
endpoints.  In particular, $A_x(0)=g_x$.
\end{lemma}

\begin{proof}
Move $x$ along a geodesic with initial velocity $\xi$ and follow the
inverse branch to the fixed endpoint $\exp_x w$.  The variational
field $Y$ is Jacobi and satisfies $Y(0)=\xi$, $Y(1)=0$.
Differentiating the first-variation formula gives
\[
 A_x(w)(\xi,\xi)=-\langle\xi,D_\tau Y(0)\rangle_x
                 =Q_{\gamma_{x,w}}(Y,Y),
\]
where the final equality is integration by parts.  There is no initial
acceleration term because the moving initial point follows a geodesic.
For any admissible $X$, put $Z=X-Y$.  The endpoint conditions and
the Jacobi equation give $Q(Y,Z)=0$, while minimality of the original
geodesic gives $Q(Z,Z)\geq0$.  Hence $Q(X,X)\geq Q(Y,Y)$.
Equality makes $Z$ a Jacobi field vanishing at both endpoints;
nonconjugacy forces $Z=0$.  At $w=0$, the minimizing field is
$Y(\tau)=(1-\tau)\xi$ along the constant geodesic, which gives
$A_x(0)(\xi,\xi)=|\xi|_x^2$.
\end{proof}

\begin{lemma}[Strict radial comparison]
\label{lem:radial-comparison}
For every $a\in G_x$ and $0<s<s'<1$,
\begin{equation}
 A_x(sa)>\frac{s}{s'}A_x(s'a)
 \label{eq:radial-comparison}
\end{equation}
as quadratic forms; that is, the difference is positive on every
nonzero vector.
\end{lemma}

\begin{proof}
Both $sa$ and $s'a$ belong to $I(x)$ by
\cref{lem:interior-characterization}.  Time rescaling in
\eqref{eq:hessian-index-minimum} gives, along $\gamma_{x,a}$,
\[
 \frac1s A_x(sa)(\xi,\xi)
 =\min_{\substack{X(0)=\xi\\X(s)=0}}Q_{\gamma_{x,a}|_{[0,s]}}(X,X).
\]
Extend the unique minimizing field by zero to $[0,s']$.
It is an admissible field for the corresponding minimum on $[0,s']$
and has the same index-form value.  If $\xi\ne0$, it cannot be that
minimum: otherwise \cref{lem:hessian-index}, with time rescaled,
would make the extension Jacobi, although it vanishes on $(s,s')$
and is nonzero at $0$.  Thus
$A_x(sa)(\xi,\xi)/s>A_x(s'a)(\xi,\xi)/s'$, proving the claim.
For $a=0$ the same argument applies, and the difference is explicitly
$(1-s/s')g_x$.
\end{proof}

\begin{lemma}[Transverse concavity]
\label{lem:transverse-concavity}
Let $b_t=b_0+t\eta$ be an affine velocity line in $T_xM$, and let
$\xi\perp\eta$.  On every interval on which $b_t\in I(x)$,
\begin{equation}
 \frac{d^2}{dt^2}A_x(b_t)(\xi,\xi)
 =-\frac23\mathfrak S_{(x,b_t)}(\xi,\eta)\leq0.
 \label{eq:transverse-concavity}
\end{equation}
Consequently, if $b=\sum_{i=1}^m\theta_i b_i$, with
$\theta_i\geq0$ and $\sum_i\theta_i=1$, the convex hull of the
$b_i$ is contained in $I(x)$, and $\langle b_i,\xi\rangle_x$
is independent of $i$, then
\[
 A_x(b)(\xi,\xi)\geq\sum_{i=1}^m\theta_iA_x(b_i)(\xi,\xi).
\]
\end{lemma}

\begin{proof}
For fixed $v\in I(x)$, the second derivative at $u=0$ of
$c(\exp_x(u\xi),\exp_x v)$ is $A_x(v)(\xi,\xi)$, because
$u\mapsto\exp_x(u\xi)$ is a geodesic.  Differentiating this identity
twice in the affine velocity direction $\eta$ and using
\eqref{eq:mtw-definition} gives the equality in
\eqref{eq:transverse-concavity}.  The inequality is exactly
\eqref{eq:weak-mtw}.  Each segment in the stated convex hull has
direction perpendicular to $\xi$, so the resulting linewise concavity
gives the finite Jensen inequality by induction.  Vanishing test
vectors or line directions cause no exception.
\end{proof}
